%% DESIGN: Palatino | Academic CV style, numbered sections | Dark Red (#8B0000) accent
\documentclass[11pt,letterpaper]{article}
\usepackage[margin=0.75in]{geometry}
\usepackage{palatino}
\usepackage[T1]{fontenc}
\usepackage{xcolor}
\usepackage{titlesec}
\usepackage{enumitem}
\usepackage{hyperref}
\usepackage{longtable}

\definecolor{darkred}{HTML}{8B0000}

\titleformat{\section}{\large\bfseries\color{darkred}}{}{0em}{}[\vspace{-0.6em}{\color{darkred}\hrule height 0.5pt}]
\titlespacing*{\section}{0pt}{12pt}{6pt}

\titleformat{\subsection}[runin]{\bfseries}{}{0em}{}[.]
\titlespacing*{\subsection}{0pt}{6pt}{4pt}

\setlist[itemize]{leftmargin=1.5em,itemsep=1pt,parsep=0pt,label=--}

\hypersetup{colorlinks=true,urlcolor=darkred,linkcolor=darkred}

\begin{document}

\begin{center}
{\LARGE\bfseries Dr.\ Rajesh Venkataraman}\\[4pt]
{\small Postdoctoral Researcher $\rightarrow$ Industry Transition}\\[4pt]
{\footnotesize Chennai, India \quad $\bullet$ \quad rajesh.v@email.com \quad $\bullet$ \quad +91-94567-12345 \quad $\bullet$ \quad \href{https://scholar.google.com}{Google Scholar} \quad $\bullet$ \quad \href{https://github.com/rajeshv}{GitHub}}
\end{center}

\section{Research Summary}
Machine learning researcher with 6 publications (h-index: 8) transitioning to industry ML/AI roles. Expertise in NLP, transformer architectures, and large-scale distributed training. Seeking applied research or senior ML engineer positions.

\section{Education}
\textbf{Ph.D.\ Computer Science} -- IIT Madras \hfill 2019 -- 2024\\
\textit{Thesis: ``Efficient Attention Mechanisms for Low-Resource Language Models''}\\
Advisor: Prof.\ Balaraman Ravindran\\[4pt]
\textbf{M.Tech Computer Science} -- IIT Madras \hfill 2017 -- 2019\\[2pt]
\textbf{B.Tech Information Technology} -- Anna University \hfill 2013 -- 2017

\section{Selected Publications}
\begin{enumerate}[leftmargin=2em,itemsep=3pt]
\item \textbf{Venkataraman, R.}, Kumar, S. ``Sparse Attention is All You Need for Low-Resource NLP.'' \textit{ACL 2024}. (Oral)
\item \textbf{Venkataraman, R.}, et al. ``Efficient Fine-tuning of Billion-Parameter Models with Adapter Fusion.'' \textit{NeurIPS 2023}.
\item \textbf{Venkataraman, R.}, Ravindran, B. ``Cross-lingual Transfer for Dravidian Languages.'' \textit{EMNLP 2022}.
\item Kumar, S., \textbf{Venkataraman, R.} ``Distillation Strategies for Multilingual Transformers.'' \textit{ICLR 2022}.
\end{enumerate}

\section{Industry Experience}
\textbf{Google Research India} \hfill Bangalore\\
\textit{Research Intern} \hfill Summer 2022
\begin{itemize}
\item Developed novel tokenization strategy for Indian languages improving BLEU scores by 4.2 points
\item Contributed to production multilingual model serving 100M+ users
\item Filed 1 patent on efficient inference for on-device language models
\end{itemize}

\textbf{Microsoft Research Lab} \hfill Hyderabad\\
\textit{Research Intern} \hfill Summer 2021
\begin{itemize}
\item Built few-shot learning pipeline for code generation reducing annotation needs by 70\%
\item Published findings at AAAI 2022 workshop on AI for Code
\end{itemize}

\section{Technical Skills}
\textbf{Languages:} Python, C++, CUDA, Julia\\
\textbf{Frameworks:} PyTorch, JAX, Hugging Face, DeepSpeed, Megatron-LM\\
\textbf{Infrastructure:} AWS (SageMaker, EC2), GCP (Vertex AI), Docker, Kubernetes\\
\textbf{Specializations:} NLP, Transformers, Distributed Training, Model Compression

\section{Awards \& Service}
\begin{itemize}
\item Best Paper Award -- ACL 2024 Student Research Workshop
\item Google PhD Fellowship Finalist (2022)
\item Reviewer: NeurIPS, ICML, ACL, EMNLP (2021--2024)
\item Teaching Assistant: Deep Learning (CS6910) -- 4 semesters
\end{itemize}

\end{document}
